% Sections 5-6 -- design ablations; limitations and robustness.  v3, 13 Oct 2026 (FINAL_SUBMISSION_REVIEW_ASTRA: compressed;
% protocol details, per-seed tables and the extensions moved to the appendix).
\section{What the learned cycle gains: design ablations}\label{sec:design}

\subsection{Deflation, output and non-linearity}\label{sec:ablation}
\begin{table}[tbp]\centering\footnotesize\setlength{\tabcolsep}{4pt}
\caption{Separately trained design variants, held-out data seed 1. Deflation, output used at inference (plain cycle or Galerkin projection), training loss and smoother class; mean energy error with each model's own output and reference/model $G$, $W$. OA-VCQ was not run on channels, where no near-kernel vectors are detected.}\label{tab:factorial}
\begin{adjustbox}{max width=\linewidth}
\begin{tabular}{llllllllll}\toprule
 & deflation & output & loss & smoothers & params & \multicolumn{2}{c}{barriers: energy, ($G$, $W$)} & \multicolumn{2}{c}{channels: energy, ($G$, $W$)} \\\midrule
OA-VC & no & plain & plain & learned & 20{,}672 & 2.33e-1 & 0.71, 0.44 & 4.08e-3 & 8.83, 0.99 \\
OA-VCZ & yes & plain & plain & learned & 20{,}672 & 7.51e-2 & 1.52, 0.75 & 3.52e-3 & 8.72, 0.99 \\
OA-VCQ & no & Galerkin & plain + proj. & learned & 20{,}672 & 1.07e-1 & 1.15, 0.65 & -- & -- \\
OA-VCB & yes & Galerkin & plain + proj. & learned & 20{,}672 & 4.36e-2 & 2.13, 0.92 & 2.55e-3 & 10.23, 1.00 \\
OA-VCL & yes & plain & plain & residual-linear & 21{,}600 & 1.51e-1 & 0.70, 0.16 & 1.11e-2 & 2.59, 0.97 \\
OA-VCBL & yes & Galerkin & plain + proj. & residual-linear & 21{,}600 & 6.76e-2 & 1.13, 0.69 & 1.01e-2 & 2.72, 0.99 \\
\bottomrule\end{tabular}
\end{adjustbox}
\end{table}

Table~\ref{tab:factorial} compares separately trained design variants on held-out seed 1; \appref{app:repeat} evaluates each
trained network with both outputs, plain and projected (Table~\ref{tab:factorial8}). Except for the residual-linear channel
control, each comparison rests on one initialisation seed per model.
\emph{Deflation.} On barriers, adding indicator deflation improves the paired accuracy by $G = 2.14$ with plain output
(OA-VC $\to$ OA-VCZ) and by $1.85$ with Galerkin output (OA-VCQ $\to$ OA-VCB). Without it no variant beats the reference by more
than $15\%$. On channels deflation is inactive.
\emph{Galerkin output and its training term.} On barriers the Galerkin output with the additional loss term adds $G = 1.63$
without deflation (OA-VC $\to$ OA-VCQ) and $1.40$ with it (OA-VCZ $\to$ OA-VCB); on channels OA-VCB is $1.17$ times more accurate
than OA-VCZ. The two components interact: without deflation, training for the projected output improves the projected output
($0.121\to0.107$) but degrades the plain one ($0.233\to0.311$); with deflation both improve ($0.049\to0.044$ projected,
$0.075\to0.066$ plain). The gains are therefore not independent and additive.
\emph{Non-linearity.} The non-linear learned smoothers outperform the tested residual-linear controls, which keep the same
operator-dependent structure but make the smoother maps linear in their right-hand side (OA-VCL is linear in $r$ for each
$A$; OA-VCBL's Galerkin output is not). On channels the residual-linear models still gain over the Chebyshev reference
($G = 2.59$ for OA-VCL, $2.72$ for OA-VCBL), and the non-linear smoothers add a further $3.36$ (plain output) and $3.76$ (Galerkin
output); on barriers the residual-linear models gain little or lose ($G = 1.13$ and $0.70$), and non-linearity adds $1.9$--$2.2$.
The product $2.72\times3.76 = 10.2$ follows algebraically from chaining paired geometric-mean ratios; it is not a causal
decomposition. Over three initialisation seeds of OA-VCBL the two factors are $2.63$--$2.72$ and $3.76$--$4.08$ (OA-VCB of the
same seed index). One OA-VCBL run (seed 2) ended with a different optimisation outcome, a $2.8$ times higher final training loss
and a plain-cycle error of $5.9\times10^{-2}$ instead of $1.07\times10^{-2}$, while its projected mean error is similar
($9.8\times10^{-3}$ against $1.01\times10^{-2}$; its paired gain is slightly lower, $2.63$): projected accuracy is stable despite
the sensitivity of the plain cycle and of the training objective.

\subsection{Operator-adapted transfers}\label{sec:transfers}
\begin{table}[tbp]\centering\footnotesize\setlength{\tabcolsep}{4pt}
\caption{Operator-adapted ($m=4$) vs geometric ($m=0$) transfers, separately trained OA-VCB models on the same data. OA-VCB error with each; paired $G$ (geometric-transfer error / adapted-transfer error) and $W$; the same ratio for the two classical references; and reference/learned $G$ for each choice of transfers.}\label{tab:transfers}
\begin{adjustbox}{max width=\linewidth}
\begin{tabular}{llllllll}\toprule
case & adapted & geometric & $G$ geom./adapt. & $W$ & references: geom./adapt. & reference/learned (adapt.) & reference/learned (geom.) \\\midrule
channels, seed 1 & 2.55e-3 & 5.02e-3 & 1.91 [1.84, 2.00] & 1.00 & 1.24 & 10.23 & 6.65 \\
channels, seed 2 & 2.79e-3 & 5.46e-3 & 1.92 [1.84, 2.00] & 1.00 & 1.24 & 9.76 & 6.30 \\
channels, seed 3 & 2.92e-3 & 5.66e-3 & 1.87 [1.80, 1.95] & 1.00 & 1.25 & 9.68 & 6.50 \\
barriers, seed 1 & 4.36e-2 & 1.13e-1 & 4.07 [3.50, 4.71] & 0.97 & 2.82 & 2.13 & 1.48 \\
\bottomrule\end{tabular}
\end{adjustbox}
\end{table}

To isolate the operator-adapted transfers, OA-VCB was retrained with $m=0$ in \eqref{eq:jacobi} (the geometric prolongation)
on the same data (Table~\ref{tab:transfers}). The adapted transfers make the learned cycle $1.9$ times more accurate on
channels, on every sample of all three seeds, and $4.1$ times more accurate on barriers ($97\%$ of samples). They help the
learned cycle more than the classical one ($1.24$ and $2.82$ for the references), and the margin of learning over the classical
reference is larger with them ($10.2$ against $6.7$ on channels, seed 1; $2.13$ against $1.48$ on barriers). Because the two
models are trained separately, the comparison measures the choice of transfers under our training protocol, including its
effect on optimisation. The adapted transfers are also more expensive: operator data $121$--$152$ instead of $11$--$14$\,ms per
sample on the training platform's CPU. Retraining with $m = 8$ keeps the margin over the correspondingly stronger reference
($G = 11.17$ on channels, $2.01$ on barriers; \appref{app:sensitivity}), so it is not an artefact of the cheaper transfers.

\section{Limitations and robustness}\label{sec:failure}

\subsection{Floating islands}\label{sec:islands}
Floating islands are present in $39$--$48\%$ of the barrier test samples, and island samples have larger mean energy errors
for every method tested (Table~\ref{tab:islands}). Indicator deflation improves the tested barrier models
(Section~\ref{sec:ablation}), while OA-VCB's remaining island errors ($0.10$--$0.14$) are still large. An island indicator is
constant on the island, whereas the low-energy mode it approximates also varies across the surrounding barrier. We therefore
asked, with a protocol fixed in writing before any number was computed (\appref{app:harmonic}), how much of OA-VCB's remaining
error lies in $A$-harmonic extensions of the indicators into a patch of $\rho = 2$ graph rings, added as a sequential
$A$-orthogonal correction of the frozen prediction; the outcome was decided on $100$ fresh coefficient fields per checkpoint.
On the 89 detector-positive fresh operators, the correction reduces OA-VCB's mean energy error from $0.1219$ to $0.1043$
($R = 14.4\%$, paired 95\% interval $9.8$--$19.5\%$; Table~\ref{tab:harmonic}) and outperforms a patch-indicator control of
equal support by a paired factor $1.13$, but it does not meet the prespecified support threshold: the outcome is ``little
support from the tested enrichment''. On a fixed 40-sample subset, extending the basis globally raises the reduction from
$10.9\%$ to $20.9\%$, so localisation matters, but substantial error remains. After the same correction is applied to both,
OA-VCB keeps its advantage over the reference on island samples ($G = 1.42$ [1.27, 1.60]). These results concern sequential
post-processing of frozen predictions; they do not assess an enlarged output space or retraining with enriched deflation.

\subsection{Dependence on the load}\label{sec:rhs}
\begin{table}[tbp]\centering\footnotesize\setlength{\tabcolsep}{4pt}
\caption{Load dependence (channels, held-out data seeds 1--3, the same 300 operators for every load; $u_c$ recomputed for each load). Pooled reference/OA-VCB $G$ [95\% CI] and $W$ for the models trained with the single load $f=1$ and for the models trained on the random-load distribution (one model per data seed; $W$: smallest of the three per-seed win fractions), and the paired ratio of the fixed-load models' error to the random-load models' (above 1 favours the random-load models). Per-seed values and mean errors: Tables~\ref{tab:rhs} and~\ref{tab:rhsrandom}.}\label{tab:loads}
\begin{adjustbox}{max width=\linewidth}
\begin{tabular}{llllll}\toprule
load & fixed-load $G$ [95\% CI] & $W$ & random-load $G$ [95\% CI] & $W$ (min) & fixed/random \\\midrule
$f=1$ & 9.89 [9.57, 10.20] & 1.00 & 6.87 [6.62, 7.13] & 1.00 & 0.70 \\
$f=-1$ & 0.21 [0.20, 0.23] & 0.00 & 6.51 [6.27, 6.75] & 1.00 & 30.8 \\
$f=3$ (homogeneity check) & 9.89 [9.57, 10.20] & 1.00 & 6.87 [6.62, 7.13] & 1.00 & 0.70 \\
random-load distribution (Sec.~\ref{sec:k0}) & 0.60 [0.50, 0.72] & 0.37 & 6.22 [5.99, 6.44] & 1.00 & 10.3 \\
smooth, both signs & 0.20 [0.19, 0.23] & 0.06 & 6.38 [6.20, 6.56] & 1.00 & 31.1 \\
localised source & 0.44 [0.38, 0.51] & 0.32 & 2.51 [2.25, 2.79] & 0.79 & 5.7 \\
\bottomrule\end{tabular}
\end{adjustbox}
\end{table}

The models of Section~\ref{sec:accuracy} were trained with the single load $f=1$ and approximate the fixed-load map
$a\mapsto u(a;f)$. Table~\ref{tab:loads} applies them and the models trained on a load distribution (Section~\ref{sec:k0}) to
other loads on the same test operators. Positive scaling ($f=3$) reproduces the relative errors of $f=1$, as
Proposition~\ref{prop:homogeneity} guarantees. The other tested load families reverse the fixed-load models' advantage: with
$f=-1$ OA-VCB is less accurate than the reference on all 300 operators ($G = 0.21$), and for smooth loads of both signs,
localised sources and the random-load distribution the reference is more accurate as well ($G = 0.20$--$0.60$). The energy
safeguard holds throughout (never worse than $u_c$ on the $1{,}800$ load--sample pairs) and shows what it promises: no damage,
not accuracy. Training on the load distribution changes this: the random-load models are more accurate than the reference by
$G = 6.2$--$6.9$ on $f=\pm1$, smooth loads and held-out random loads, with $W = 1$ on every seed, at the price of a $1.44$-fold
loss on $f=1$ against fixed-load training. On pure localised sources the pooled gain is $2.51$, but the per-seed win fractions
are $0.89$, $0.79$ and $0.83$, so the criterion fixed before this test ($G\ge2$ and $W\ge0.8$ per seed) is missed on seed 2, by one
sample. These tests support broader forcing coverage after distributional training; they do not establish arbitrary-load
accuracy (load definitions and per-seed values in \appref{app:loads}).

\subsection{Further settings}\label{sec:k0}\label{sec:l4}\label{sec:aniso}
\begin{table}[tbp]\centering\footnotesize\setlength{\tabcolsep}{4pt}
\caption{Further settings: mean energy error of OA-VCB and of its classical reference, paired reference/OA-VCB $G$ [95\% CI] and $W$, and the per-sample number of two-grid PCG iterations, from the same starting point, at which PCG first reaches OA-VCB's energy error (median, range; computed iteration histories: 12 for $k_0=0$, 16 otherwise).}\label{tab:settings}
\begin{adjustbox}{max width=\linewidth}
\begin{tabular}{llllll}\toprule
setting & OA-VCB & reference & $G$ [95\% CI] & $W$ & PCG its \\\midrule
$k_0=0$ (no external coarse solution), seed 1 & 7.62e-3 & 8.52e-2 & 11.39 [10.95, 11.83] & 1.00 & 5 (4--7) \\
$k_0=0$ (no external coarse solution), seed 2 & 7.58e-3 & 8.04e-2 & 10.91 [10.35, 11.44] & 1.00 & 5 (4--6) \\
$k_0=0$ (no external coarse solution), seed 3 & 8.13e-3 & 8.55e-2 & 10.82 [10.40, 11.26] & 1.00 & 5 (5--7) \\
finer hierarchy (L4, $k_0=3$), seed 1 & 2.21e-3 & 2.01e-2 & 9.72 [9.04, 10.33] & 1.00 & 6 (4--8) \\
polygonal domain, seed 1 & 3.54e-3 & 3.30e-2 & 9.39 [8.89, 9.86] & 1.00 & 6 (5--8) \\
random loads (in distribution), seed 1 & 4.50e-3 & 2.65e-2 & 6.27 [5.89, 6.65] & 1.00 & 5 (4--7) \\
random loads (in distribution), seed 2 & 4.26e-3 & 2.51e-2 & 6.05 [5.72, 6.36] & 1.00 & 5 (4--7) \\
random loads (in distribution), seed 3 & 4.54e-3 & 2.61e-2 & 6.34 [5.90, 6.77] & 1.00 & 6 (4--7) \\
anisotropic tensor coefficient, seed 0 (dev.) & 6.21e-3 & 2.10e-2 & 3.50 [3.33, 3.68] & 1.00 & -- \\
anisotropic tensor coefficient, seed 1 & 6.01e-3 & 2.04e-2 & 3.51 [3.33, 3.69] & 1.00 & -- \\
anisotropic tensor coefficient, seed 2 & 6.17e-3 & 2.12e-2 & 3.58 [3.40, 3.77] & 1.00 & -- \\
\bottomrule\end{tabular}
\end{adjustbox}
\end{table}

Table~\ref{tab:settings} summarises separately trained variations of the main configuration (details in
\appref{app:settings}). Without an externally supplied coarse solution ($k_0=0$; the cycle still contains its exact 28-node
coarsest solve) OA-VCB keeps $G = 10.8$--$11.4$ with $W = 1$, and the energy safeguard then only bounds the damage relative to
$u=0$. A separately trained finer hierarchy ($L=4$, 14{,}398 unknowns) gives $G = 9.72$, a polygonal domain $9.39$, an anisotropic
tensor coefficient $3.5$--$3.6$, and training and testing on the random-load distribution $6.05$--$6.34$, all with $W=1$.
\emph{Refinement without retraining.} The L4-trained model applied unchanged to the next refined mesh (L5, 58{,}494 unknowns,
100 test samples) keeps $G = 8.12$ [7.50, 8.67] with $W = 1$, a per-sample gain $0.84$ [0.76, 0.91] times that at L4 on the same
coefficient draws. At L6 (235{,}774 unknowns) the evaluation was stopped by the platform's session limit after 21 of 24 planned
samples; on these completed cases $G = 4.18$ [3.43, 4.88] with $W = 1$, about half the L5 gain on the same draws. This partial
evaluation shows a retained but shrinking margin; it does not establish mesh independence, and the coarse level, the bottom
level of the cycle and the distance from the training mesh change together. A model trained at L3 does not keep its margin at
L4 ($G = 1.95$ on 20 samples). \emph{Plane-strain elasticity} (one data set, one training run; preliminary): without deflation and without
tuning, OA-VCB is more accurate than its Chebyshev reference on every test sample ($G = 2.32$; \appref{app:settings}).
